\documentclass[11pt]{article}
\usepackage[margin=1in]{geometry}
\usepackage{amsmath,amssymb,amsthm}
\usepackage{hyperref}
\usepackage{enumitem}

\newtheorem{theorem}{Theorem}
\newtheorem{lemma}[theorem]{Lemma}
\newtheorem{corollary}[theorem]{Corollary}
\theoremstyle{definition}
\newtheorem{definition}[theorem]{Definition}
\theoremstyle{remark}
\newtheorem{remark}[theorem]{Remark}

\usepackage{xurl}
\let\oldus\_
\renewcommand{\_}{\oldus\allowbreak}
\newcommand{\M}{\mathcal{M}(\mathbb{C})}

\title{A disproof of Conjecture 00000002251\\[2pt]
\large The kernel of the difference operator on meromorphic functions is not one-dimensional}
\date{}

\begin{document}
\sloppy
\maketitle

\section{The conjecture}

Conjecture 00000002251 reads (English version):

\begin{quote}
\emph{Definition:} The difference operator $\Delta_c f = f(z+c) - f(z)$.
\emph{Conjecture:} The kernel of $\Delta_c$ has dimension 1, and the spectrum of $\Delta_c$ on
meromorphic function spaces (differences of difference polynomials) is explicit of the
$\{2\sin(kc/2)\}$ type (difference spectrum).
\end{quote}

The Chinese version says the same: the kernel of $\Delta_c$ has dimension $1$, and the spectrum of
$\Delta_c$ on the meromorphic function space (\emph{ya chun han shu kong jian}) is explicitly of
$\{2\sin(kc/2)\}$ type. The statement is a conjunction. Its first clause (kernel of dimension $1$) is
unhedged; the second clause is hedged by the word ``type''. We refute the first clause, which refutes the
conjunction, and we also compute the spectrum exactly.

\section{Reading and conventions}

\begin{definition}
$\M$ is the set of meromorphic functions $f:\mathbb{C}\to\mathbb{C}$ (meromorphic at every point of
$\mathbb{C}$), a complex vector space under pointwise operations. For $c\in\mathbb{C}$,
$\Delta_c:\M\to\M$ is $(\Delta_c f)(z) = f(z+c)-f(z)$; it maps $\M$ into $\M$ because translates, sums and
scalar multiples of meromorphic functions are meromorphic, and it is $\mathbb{C}$-linear.
\end{definition}

\begin{itemize}[nosep]
\item \textbf{Domain.} We read $\Delta_c$ on the space the text names, ``meromorphic function spaces'',
  with no growth, periodicity or polynomial restriction (Section~\ref{sec:readings} lists restricted
  readings, which we do \emph{not} refute).
\item \textbf{Values at poles.} Following Mathlib, a meromorphic function is an honest function
  $\mathbb{C}\to\mathbb{C}$ with some value at each pole. Lemma~\ref{lem:germ} shows that the witnesses
  stay independent if functions that agree off a discrete set are identified, so the result does not
  depend on this convention.
\item \textbf{Dimension} is the rank of a vector space (a cardinal, \texttt{Module.rank}); ``dimension 1''
  means rank $=1$. \textbf{Spectrum} is the algebraic spectrum in $\mathrm{End}_{\mathbb{C}}(\M)$:
  $\mu\in\sigma(\Delta_c)$ iff $\mu\cdot\mathrm{id}-\Delta_c$ is not bijective.
\item \textbf{The shift $c$} ranges over all of $\mathbb{C}$; the kernel result holds for every $c$,
  including the degenerate $c=0$. The spectrum result is stated for $c\neq 0$.
\end{itemize}

\section{The kernel has dimension at least two}

\begin{lemma}\label{lem:eigen}
For $b,c\in\mathbb{C}$ let $e_b(z)=e^{bz}$. Then $\Delta_c e_b = (e^{bc}-1)\,e_b$.
\end{lemma}
\begin{proof}
$e^{b(z+c)}-e^{bz} = e^{bz}e^{bc}-e^{bz}$.
\end{proof}

\begin{theorem}\label{thm:main}
For every $c\in\mathbb{C}$, $\operatorname{rank}_{\mathbb{C}}\ker\Delta_c > 1$. In particular
$\ker \Delta_c$ never has dimension $1$.
\end{theorem}
\begin{proof}
\emph{Case $c\neq 0$.} Let $\omega = 2\pi i/c$, $u=1$ and $v=e_\omega$. Then $\Delta_c u = 0$ and, by
Lemma~\ref{lem:eigen}, $\Delta_c v = (e^{2\pi i}-1)v = 0$. If $su+tv = s'u+t'v$, evaluate at $z=0$ and at
$z=c/2$, where $v(0)=1$ and $v(c/2)=e^{\pi i}=-1$: $s+t=s'+t'$ and $s-t=s'-t'$, hence $s=s'$, $t=t'$. So
$u,v$ are linearly independent elements of $\ker\Delta_c$.

\emph{Case $c=0$.} $\Delta_0=0$, so $u=1$ and $v(z)=z$ lie in the kernel; evaluating at $z=0$ and $z=1$
gives independence as above.

In both cases $\ker\Delta_c$ contains a linearly independent pair, so its rank is at least $2>1$.
\end{proof}

\begin{lemma}\label{lem:germ}
Let $c\neq0$ and $s,t\in\mathbb{C}$. If $s+t\,e^{2\pi i z/c}=0$ for all $z$ outside a discrete subset
of $\mathbb{C}$ (i.e.\ on a set in the codiscrete filter), then $s=t=0$.
\end{lemma}
\begin{proof}
The function $F(z)=s+te^{2\pi iz/c}$ is continuous and vanishes on a punctured neighbourhood of every
point, hence everywhere. Now use $z=0$ and $z=c/2$ as before.
\end{proof}

\begin{remark}
Both witnesses for $c\neq 0$ are entire, and $|e^{bz}|=e^{\operatorname{Re}(bz)}\le e^{|b||z|}$. Hence the
same pair lies in the kernel of $\Delta_c$ on entire functions and on any class of meromorphic functions
that contains entire functions of exponential type, so those readings are refuted as well.
\end{remark}

\section{The spectrum}

\begin{theorem}\label{thm:spec}
Let $c\neq0$. Every $\mu\neq-1$ is an eigenvalue of $\Delta_c$ on $\M$, and $-1\notin\sigma(\Delta_c)$.
Hence $\sigma(\Delta_c)=\mathbb{C}\setminus\{-1\}$, which is uncountable; in particular
$\sigma(\Delta_c)\ne\{2\sin(kc/2):k\in\mathbb{Z}\}$, and it is not the image of any countable index set.
\end{theorem}
\begin{proof}
For $\mu\neq-1$ put $b=\log(1+\mu)/c$ (any branch; $1+\mu\neq0$). Then $e^{bc}=1+\mu$, and by
Lemma~\ref{lem:eigen} $\Delta_c e_b=\mu e_b$ with $e_b(0)=1\neq0$. For $\mu=-1$:
$(-\mathrm{id}-\Delta_c)f(z) = -f(z+c)$, which is a bijection of $\M$ with inverse $g\mapsto
-g(\cdot-c)$. Finally $\mathbb{C}=(\mathbb{C}\setminus\{-1\})\cup\{-1\}$ is uncountable, while the image of
a countable set is countable.
\end{proof}

Note that $0\in\sigma(\Delta_c)$: it is the eigenvalue of the (non-trivial) kernel.

\section{Relation to the earlier PR \#270}

PR \#270 (orionsheep, merged 2026-10-03) used the same reading and the same witness
$e^{2\pi iz/c}$. It was removed in the re-audit commit \texttt{541cf4fb}, whose note reads
``\#2251 (4-point cyclic shadow; meromorphic kernel prose-only)''. Its Lean file defined a function
\texttt{unit4 : Nat -> Int x Int} cycling through four pairs labelled $1,i,-1,-i$ and proved
\texttt{unit4 (z+4) = unit4 z} by \texttt{rfl}, plus the distinctness of four integer pairs. No complex
function, no meromorphic function, no operator $\Delta_c$ and no kernel or dimension appeared in Lean; the
theorem that the kernel has dimension $\ge2$ was prose only, so the Lean certified nothing about the
conjecture's objects.

PR \#270 also stated that the eigenvalues $e^{\lambda c}-1$ of the eigenfunctions $e^{\lambda z}$ ``sweep
$\mathbb{C}\setminus\{0\}$''. That is wrong: as $\lambda$ ranges over $\mathbb{C}$ (and $c\ne0$),
$e^{\lambda c}$ takes every nonzero value, so $e^{\lambda c}-1$ takes exactly the values in
$\mathbb{C}\setminus\{-1\}$. The value $0$ \emph{is} attained (it is the eigenvalue of the non-trivial
kernel, contradicting PR \#270's own kernel claim), and $-1$ is the omitted value
(Theorem~\ref{thm:spec}). The present submission corrects both defects: $\Delta_c$ is a
$\mathbb{C}$-linear endomorphism of the genuine space of meromorphic functions $\mathbb{C}\to\mathbb{C}$,
and the rank of its kernel and its spectrum are theorems about that object.

\section{Readings that are not refuted}\label{sec:readings}

The kernel clause can hold if $\Delta_c$ is restricted to a smaller space that excludes non-constant
periodic functions. We do \emph{not} refute such readings, for example: $\Delta_c$ on polynomials or
rational functions (for $c\neq 0$ a $c$-periodic one is constant: if $z_0$ is not a pole, then
$f-f(z_0)$ vanishes at $z_0+nc$ for all $n\in\mathbb{N}$); on meromorphic functions of order $<1$; or on $2\pi$-periodic trigonometric polynomials
$\sum a_k e^{ikz}$ with $c/2\pi$ irrational. In the last setting
$\Delta_c e^{ikz}=(e^{ikc}-1)e^{ikz}=2i\,e^{ikc/2}\sin(kc/2)\,e^{ikz}$, which is presumably where the
$\{2\sin(kc/2)\}$ set comes from.

We take the literal reading because the text itself names the domain: ``the spectrum of $\Delta_c$ on
meromorphic function spaces'', with ``difference polynomials'' (polynomial expressions in shifts
$f(z+c_j)$) as the objects $\Delta_c$ acts on. No growth bound, periodicity or polynomial restriction is
stated, and the kernel clause carries no hedge. On that space the kernel of $\Delta_c$ is the space of all
$c$-periodic meromorphic functions, which has dimension at least $2$ for every $c$.

\section{Lean 4 formalization}

The file \texttt{lean/Conjecture2251/Basic.lean} (221 lines, only \texttt{import Mathlib}; Lean
\texttt{v4.33.1}, Mathlib \texttt{v4.33.1}) defines
\begin{itemize}[nosep]
\item \texttt{MeroFun : Submodule C (C -> C)}, carrier $\{f \mid \texttt{Meromorphic } f\}$ (Mathlib's
  \texttt{Meromorphic}: meromorphic at every point);
\item \texttt{diffOp c : Module.End C MeroFun}, $f\mapsto$ \texttt{fwdDiff c f}, i.e.\
  $z\mapsto f(z+c)-f(z)$ (Mathlib's forward difference), with closure from
  \texttt{Meromorphic.meromorphic\_fun\_comp\_add\_const\_iff\_meromorphic}.
\end{itemize}
Theorems (all in namespace \texttt{Tlmc2251}):
\begin{itemize}[nosep]
\item \texttt{diffOp\_expMF}: $\Delta_c e_b=(e^{bc}-1)e_b$ (Lemma~\ref{lem:eigen});
  \texttt{expMF\_differentiable}, \texttt{expMF\_growth}: $e_b$ is entire and $\|e^{bz}\|\le
  e^{\|b\|\|z\|}$;
\item \texttt{kernel\_pair}: for every $c$, two elements of $\ker\Delta_c$ with
  $su+tv=s'u+t'v\Rightarrow s=s',t=t'$;
\item \texttt{rank\_ker\_diffOp\_gt\_one}: \texttt{1 < Module.rank C (LinearMap.ker (diffOp c))} for all
  \texttt{c} (Theorem~\ref{thm:main}); \texttt{kernel\_not\_dim\_one}: the rank is $\ne 1$;
\item \texttt{indep\_mod\_codiscrete}: Lemma~\ref{lem:germ}, with the filter \texttt{codiscrete C};
\item \texttt{hasEigenvalue\_diffOp}, \texttt{neg\_one\_not\_mem\_spectrum}, \texttt{spectrum\_diffOp}
  (\texttt{spectrum C (diffOp c) = \{-1\}\textasciicircum c} for $c\ne0$),
  \texttt{spectrum\_not\_countable}, \texttt{spectrum\_ne\_sine\_set} (Theorem~\ref{thm:spec});
\item \texttt{conjecture\_2251\_refuted}: the conjunction of \texttt{rank\_ker\_diffOp\_gt\_one} for all
  $c$ and, for $c\ne0$, \texttt{spectrum\_diffOp} and \texttt{spectrum\_not\_countable}.
\end{itemize}
Here \texttt{C} stands for $\mathbb{C}$. \texttt{lake build} succeeds, and \texttt{\#print axioms} for
\texttt{conjecture\_2251\_refuted}, \texttt{kernel\_not\_dim\_one}, \texttt{indep\_mod\_codiscrete},
\texttt{spectrum\_ne\_sine\_set} and \texttt{expMF\_growth} reports only \texttt{propext},
\texttt{Classical.choice} and \texttt{Quot.sound}. There is no \texttt{sorry}, no \texttt{native\_decide}
and no added axiom.

\end{document}
